\section*{Chapitre \ref{ch:fm:the-economics-of-a-trading-firm} --- L'économie d'une société de trading}

\begin{solution}{exo:fm:the-economics-of-a-trading-firm:1}
$2\,436.7+617.0+508.8-769.8-647.4=2\,145.3$ millions de dollars. \omterm{def:fm:the-economics-of-a-trading-firm:ntr}{Coûts liés au volume} $769.8+647.4=1\,417.2$ face au revenu avant autres produits
$2\,145.3+1\,417.2=3\,562.5$~: 39.8~\%.
\end{solution}

\begin{solution}{exo:fm:the-economics-of-a-trading-firm:2}
$675/1\,398=48.3~\%$ en 2025 et $678/1\,696=40.0~\%$ en 2022. La rémunération a baissé de 3 millions de dollars alors que le revenu a baissé de 298 millions de dollars~: l'essentiel de la rémunération ne suit pas
le revenu d'une année sur l'autre, si bien que le ratio évolue à rebours du revenu.
\end{solution}

\begin{solution}{exo:fm:the-economics-of-a-trading-firm:3}
$13\,117/79\,748=16.4~\%$, comme publié.
\end{solution}

\begin{solution}{exo:fm:the-economics-of-a-trading-firm:4}
Revenu net 75. Rémunération variable~: $0.7\times75-50=2.5$. Rémunération entièrement fixe~: $75-80=-5$, soit une perte.
\end{solution}

\begin{solution}{exo:fm:the-economics-of-a-trading-firm:5}
Revenu net $+133.0~\%$, résultat opérationnel $+607.5~\%$, ratio 4.57. Le levier opérationnel au revenu de 2019 vaut $1/x^\ast=1/0.229=4.37$ avec la
pente ajustée de la rémunération (4.59 avec une rémunération fixe)~: le ratio d'arc sur un doublement est proche de l'élasticité ponctuelle, parce que les coûts ont à peine bougé.
\end{solution}

\begin{solution}{exo:fm:the-economics-of-a-trading-firm:6}
2025~: rémunération de 528.1 millions de dollars contre 438.0 millions de dollars sur la ligne de revenu. Vérifier si la rémunération a commencé à suivre le revenu (une part variable plus importante dans une année forte)
et si l'entreprise a grandi~: ses effectifs sont passés d'environ 969 à environ 1\,027 entre les deux dépôts.
\end{solution}

\begin{solution}{exo:fm:the-economics-of-a-trading-firm:7}
Sans 2020, la pente vaut 0.13 avec une erreur type de 0.13, contre $0.29\pm0.094$~: une seule année à faible revenu et à faible rémunération a porté l'estimation. La
pente n'est qu'une première lecture, et la répartition propre à l'entreprise entre rémunération fixe et variable est la meilleure source.
\end{solution}

\begin{solution}{exo:fm:the-economics-of-a-trading-firm:8}
Les dénominateurs diffèrent~: les fonds propres d'un teneur de marché sont ce que ses propriétaires ont laissé dans l'entreprise, qui finance un bilan bien plus grand grâce aux courtiers
et aux prêteurs~; les fonds propres d'un segment bancaire sont une allocation fixée par les règles de fonds propres du groupe. Les années diffèrent aussi~: le ROE du teneur de marché était de 17.3~\% en
2023. Les \omterm{def:fm:the-economics-of-a-trading-firm:roe}{rentabilités des fonds propres} ne permettent de comparer des activités que si les fonds propres sont mesurés de la même façon et sur un cycle complet.
\end{solution}

\begin{solution}{pb:fm:the-economics-of-a-trading-firm:1}
\begin{enumerate}
\item Le revenu moins les coûts qui croissent avec l'activité réalisée~; le financement des positions et des marges croît avec les positions, si bien qu'il s'agit d'un coût lié au volume.
\item 2\,145.3 millions de dollars, égal au \omterm{def:fm:the-economics-of-a-trading-firm:ntr}{revenu net de trading} ajusté de l'entreprise.
\item Teneur de marché 1\,205.7 millions de dollars, 56.2~\%~; société de gestion 428 millions de dollars, 32.3~\%~; segment bancaire 17\,574 millions de dollars, 42.8~\%.
\item Teneur de marché 2.09 millions de dollars et 0.51 million de dollars par tête~; société de gestion 0.77 million de dollars et 0.39 million de dollars.
\item Le nommer et le définir, présenter avec lui le chiffre réglementaire le plus proche (GAAP), et réconcilier quantitativement les deux.
\item $\Pi=(1-b)\mathrm{Rev}-C^{\mathrm{fix}}$.
\item $x^\ast=1-C^{\mathrm{fix}}/((1-b)\mathrm{Rev})=\Pi/((1-b)\mathrm{Rev})$~; le levier est $(1-b)\mathrm{Rev}/\Pi=1/x^\ast$.
\item Teneur de marché $0.048\pm0.045$~; société de gestion $0.29\pm0.094$~; segment bancaire 0.20 sur trois ans. Seule celle de la société de gestion se distingue
de zéro, et même celle-ci tombe à 0.13 sans 2020.
\item 836.8 millions de dollars, 510.8 millions de dollars et 15\,399.9 millions de dollars.
\item Une partie de sa rémunération est fixe~: quand le revenu baisse, la rémunération baisse moins.
\item Rémunération variable~: 59.0~\%, 45.6~\%, 53.3~\%. Rémunération maintenue fixe~: 56.2~\%, 32.3~\%, 42.8~\%.
\item Rémunération variable~: 1.69, 2.19, 1.88. Rémunération maintenue fixe~: 1.78, 3.10, 2.34.
\item 49.2~\%, 34.2~\% et 43.7~\% du résultat opérationnel de 2025.
\item Les trois résistent à une baisse d'un tiers avec une rémunération variable~; la société de gestion seulement parce que sa rémunération varie (baisse au point mort de 32.3~\% avec une rémunération fixe).
\item 21.8~\%~; baisse au point mort de 22.9~\% au revenu de 2019.
\item Teneur de marché 52 millions de dollars par point d'indice (erreur type 33), corrélation 0.58~; société de gestion $-13$ millions de dollars (23), corrélation $-0.27$.
\item 52.7~\% en 2025 et 17.3~\% en 2023.
\item Les fonds propres du teneur de marché sont le capital de ses propriétaires, qui soutient un bilan financé plus important~; ceux du segment bancaire sont une allocation des fonds propres du groupe
selon les règles de fonds propres.
\item Le tableau du chapitre~: 59.0~\%/1.69 et 56.2~\%/1.78~; 45.6~\%/2.19 et 32.3~\%/3.10~; 53.3~\%/1.88 et 42.8~\%/2.34.
\item Rendre variable une plus grande part de la rémunération, pour qu'elle suive le revenu, et réduire les coûts fixes~: les deux augmentent la baisse au point mort. Une marge étroite ne peut pas être corrigée
en espérant une année volatile.
\end{enumerate}
\end{solution}

\begin{solution}{iq:fm:the-economics-of-a-trading-firm:1}
La rémunération rapportée au revenu net sur la période. Une partie de la rémunération est fixe, si bien que lorsque le revenu baisse plus vite que la rémunération, le ratio augmente, même si la rémunération baisse.
\iqlookfor{la définition et l'argument de la part fixe.}
\end{solution}

\begin{solution}{iq:fm:the-economics-of-a-trading-firm:2}
Le volume augmente avec la volatilité, et le taux de capture par unité (fourchettes plus larges) augmente aussi~; les deux se multiplient pour former le revenu.
\iqlookfor{le revenu comme produit du volume par la capture, et les deux canaux.}
\end{solution}

\begin{solution}{iq:fm:the-economics-of-a-trading-firm:3}
L'entreprise à salaires fixes~: à coûts totaux égaux, son levier opérationnel $\mathrm{Rev}/\Pi$ dépasse $(1-b)\mathrm{Rev}/\Pi$, et sa baisse au point mort est
plus petite. L'entreprise à primes résiste à la baisse plus forte.
\iqlookfor{la \cref{prop:fm:the-economics-of-a-trading-firm:fall} et la relation réciproque.}
\end{solution}

\begin{solution}{iq:fm:the-economics-of-a-trading-firm:4}
Revenu 60~: $0.8\times60-50=-2$, soit une perte~; le résultat avant la baisse était de 30.
\iqlookfor{les coûts variables baissent avec le revenu, les coûts fixes non.}
\end{solution}

\begin{solution}{iq:fm:the-economics-of-a-trading-firm:5}
Les fonds propres sont mesurés différemment (capital des propriétaires face à une part attribuée des fonds propres du groupe)~; le bilan est financé autrement~;
l'année et le cycle diffèrent~; les règles de fonds propres diffèrent par nature (chapitre 2 face au chapitre 5).
\iqlookfor{le dénominateur, et non le numérateur.}
\end{solution}

\begin{solution}{iq:fm:the-economics-of-a-trading-firm:6}
Quatre degrés de liberté~: à environ trois erreurs types de zéro, mais une seule année peut le faire beaucoup bouger (en retirant 2020, il est divisé par deux ici). La croissance des effectifs, une
acquisition ou un changement dans la composition des commissions peut faire augmenter ensemble la rémunération et le revenu. Vérifier la répartition propre à l'entreprise entre rémunération fixe et variable.
\iqlookfor{degrés de liberté, fragilité au retrait d'une année, et facteurs confondants.}
\end{solution}
